% TODO translate
% TODO TikZ
\mysection[12ビット値を配列にパッキング]{ビット演算を使った12ビット値の配列パッキング（x64、ARM/ARM64、MIPS）}

（この部分は2015年9月4日に私のブログで最初に登場しました。）

\subsection{はじめに}

ファイルアロケーションテーブル（FAT）は広く普及したファイルシステムでした。
信じがたいことですが、おそらく単純さと互換性の理由で、フラッシュドライブでまだ使用されています。
FATテーブル自体は要素の配列で、各要素はファイルの次のクラスタ番号を指します
（FATはディスク全体に散らばったファイルをサポートします）。
これは、各要素の最大値がディスク上のクラスタの最大数であることを意味します。
MS-DOS時代、ほとんどのハードディスクはFAT16ファイルシステムを持っていました。クラスタ番号が16ビット値にパックできたからです。
その後ハードディスクが安くなり、FAT32が登場しました。ここではクラスタ番号に32ビット値が割り当てられました。
しかし、フロッピーディスケットがそれほど安くなく、多くのスペースがなかった時代もありました。そのため、すべてのファイルシステム構造をできるだけきつくパックする理由でFAT12が使用されました。

したがって、FAT12ファイルシステムのFATテーブルは、連続する2つの12ビット要素が3バイト（トリプレット）に格納される配列です。
6つの12ビット値（AAA、BBB、CCC、DDD、EEE、FFF）が9バイトにどのようにパックされるかを示します：

\begin{lstlisting}
 +0 +1 +2 +3 +4 +5 +6 +7 +8
|AA|AB|BB|CC|CD|DD|EE|EF|FF|...
\end{lstlisting}

値を配列にプッシュして取り出すことは、ビット操作演算（C/C++と低レベルマシンコードの両方で）の良い例になるので、ここではFAT12を例として使用します。

\subsection{データ構造}

各バイトトリプレットが2つの12ビット値を格納することをすぐに観察できます：最初のものは左側に、2番目のものは右側にあります：

\begin{lstlisting}
 +0 +1 +2
|11|12|22|...
\end{lstlisting}

ニブル（4ビットチャンク）を次の方法でパックします（1 - 最上位ニブル、3 - 最下位）：

（偶数）

\begin{lstlisting}
 +0 +1 +2
|12|3.|..|...
\end{lstlisting}

（奇数）

\begin{lstlisting}
 +0 +1 +2
|..|.1|23|...
\end{lstlisting}

\subsection{アルゴリズム}

アルゴリズムは次のようになります：要素のインデックスが偶数の場合、左側に配置し、インデックスが奇数の場合、右側に配置します。
中間バイト：要素のインデックスが偶数の場合、その部分を上位4ビットに配置し、奇数の場合、その部分を下位4ビットに配置します。
しかし最初に、正しいトリプレットを見つけます。これは簡単です：$\frac{index}{2}$。
バイト配列内の正しいバイトのアドレスを見つけることも簡単です：$\frac{index}{2} \cdot 3$または$index \cdot \frac{3}{2}$または単に$index \cdot 1.5$。

配列から値を取得：インデックスが偶数の場合、最左端と中間バイトを取得してその部分を結合します。
インデックスが奇数の場合、中間と最右端バイトを取得してそれらを結合します。
中間バイトで不要なビットを分離することを忘れないでください。

値のプッシュはほぼ同じですが、中間バイトの他の\IT{他者}のビットを上書きしないよう注意し、\IT{あなたの}ビットのみを修正します。

\subsection{C/C++コード}

\lstinputlisting[style=customc]{advanced/380_FAT12/FAT12.c}

テスト中、すべての12ビット要素は0..0xFFF範囲の値で埋められます。
以下はすべてのトリプレットのダンプで、各行には3バイトがあります：

\begin{lstlisting}
0x000001
0x002003
0x004005
0x006007
0x008009
0x00A00B
0x00C00D
0x00E00F
0x010011
0x012013
0x014015

...

0xFECFED
0xFEEFEF
0xFF0FF1
0xFF2FF3
0xFF4FF5
0xFF6FF7
0xFF8FF9
0xFFAFFB
0xFFCFFD
0xFFEFFF
\end{lstlisting}

また、512/2*3で始まる300バイト（または100トリプレット）のGDBバイトレベル出力があります。つまり、512番目の要素（0x200）が始まるアドレスです。
トリプレットを明示的に示すために、テキストエディタで角括弧を追加しました。
最後の要素が終了し、次の要素が始まる中間バイトに注目してください。
言い換えると、各中間バイトには偶数要素の最下位4ビットと奇数要素の最上位4ビットがあります。

\lstinputlisting{advanced/380_FAT12/gdb.txt}

\subsection{仕組み}

配列をグローバルバッファにして、よりシンプルなアクセスを可能にします。

\subsubsection{ゲッター}

まず、配列から値を取得する関数から始めましょう。これがより簡単だからです。

トリプレット番号を見つける方法は、入力インデックスを2で割るだけですが、1ビット右にシフトするだけでこれを行うことができます。
これは$2^n$の形の数による除算/乗算の非常に一般的な方法です。

どのように動作するかをデモンストレーションできます。123を10で割りたいとします。
最後の桁（3、これは除算の余り）を削除すると、12が残ります。
2による除算は最下位ビットを削除するだけです。削除は右シフトで行うことができます。

次に、関数は入力インデックスが偶数（12ビット値が左に配置）か奇数（右に配置）かを決定する必要があります。
これを行う最も簡単な方法は、最下位ビット（x\&1）を分離することです。それが0の場合、値は偶数で、そうでなければ奇数です。

この事実は簡単に説明できます。最下位ビットを見てください：

\begin{lstlisting}
decimal binary even/odd

0       0000   even
1       0001   odd
2       0010   even
3       0011   odd
4       0100   even
5       0101   odd
6       0110   even
7       0111   odd
8       1000   even
9       1001   odd
10      1010   even
11      1011   odd
12      1100   even
13      1101   odd
14      1110   even
15      1111   odd
...
\end{lstlisting}

ゼロも偶数です。2の補数システムでは、-1と1の間にあります。

数学好きには、偶数または奇数の符号も2による除算の余りであるとも言えます。
数を2で割ることは単に最後のビットを削除することで、これは除算の余りです。
まあ、ここではシフトを行う必要はありません。最下位ビットを分離するだけです。

要素が奇数の場合、中間と右バイト（\IT{array[array\_idx+1]}と\IT{array[array\_idx+2]}）を取ります。
中間バイトの最下位4ビットが分離されます。
右バイトは全体として取られます。
次に、これらの部分を12ビット値に結合する必要があります。
これを行うために、中間バイトから4ビットを8ビット左にシフトします。これらの4ビットがバイトの最上位8ビットの直後に割り当てられるようにします。
次に、OR演算を使用して、これらの部分を追加します。

要素が偶数の場合、12ビット値の上位8ビットは左バイトにあり、最下位4ビットは中間バイトの上位4ビットにあります。
中間バイトで上位4ビットを4ビット右にシフトしてから、AND演算を適用して、何も残っていないことを確認します。
また、左バイトを取り、その値を4ビット左にシフトします。これは11番目から4番目（0から始まる、包括的）のビットを持っているからです。
OR演算を使用して、これら2つの部分を結合します。

\subsubsection{セッター}

セッターは同じ方法でトリプレットのアドレスを計算します。
また、左/右バイトでも同じ方法で動作します。
しかし、中間バイトに書き込むだけでは正しくありません。書き込み操作が他の要素に関連する情報を破壊するからです。
したがって、一般的な方法は、バイトをロードし、書き込むビットを削除し、そこに書き込みますが、他の部分はそのままにしておくことです。
AND演算（C/C++で\TT{\&}）を使用して、\IT{我々の}部分以外のすべてを削除します。
OR演算（C/C++で\TT{|}）を使用して、中間バイトを更新します。

\subsection{Linux x64用最適化GCC 4.8.2}

Linux x64用最適化GCC 4.8.2が何をするかを見てみましょう。
コメントを追加しました。
時々読者は命令の順序が論理的でないことで混乱します。
これは問題ありません。最適化コンパイラはCPUのアウトオブオーダー実行メカニズムを考慮に入れるため、時々、交換された命令の方が性能が良いのです。

\subsubsection{ゲッター}

\lstinputlisting[style=customasmx86]{advanced/380_FAT12/GCC_getter.asm}

\subsubsection{セッター}

\lstinputlisting[style=customasmx86]{advanced/380_FAT12/GCC_setter.asm}

\subsubsection{その他のコメント}

Linux x64のすべてのアドレスは64ビットなので、ポインタ演算中、すべての値も64ビットでなければなりません。
配列内のオフセットを計算するコードは32ビット値で動作します（入力\IT{idx}引数は\IT{int}型で、幅は32ビット）。
したがって、これらの値は実際のメモリロード/ストア前に64ビットアドレスに変換する必要があります。
そのため、変換に使用される符号拡張命令（\INS{CDQE}、\INS{MOVSX}など）がたくさんあります。
なぜ符号を拡張するのでしょうか？C/C++標準により、ポインタ演算は負の値で動作できます
（\IT{array[-123]}のような負のインデックスを使用して配列にアクセスすることが可能です。\myref{negative_array_indices}参照）。
GCCコンパイラはすべてのインデックスが常に正であることを確信できないため、符号拡張命令を追加します。

\subsection{最適化Keil 5.05（Thumbモード）}

\subsubsection{ゲッター}

以下のコードは、関数のエピローグで最終的なOR演算を持っています。
実際、これは両方のブランチの最後で実行されるため、スペースを節約することが可能です。

\lstinputlisting[style=customasmARM]{advanced/380_FAT12/Keil_thumb_O3_getter.asm}

少なくとも冗長性があります：\IT{idx*1.5}が2回計算されます。
演習として、読者はアセンブリ関数をより短くするように書き直すことを試すことができます。テストを忘れないでください！

もう一つ言及すべきことは、16ビットThumb命令で大きな定数を生成するのが難しいため、Keilコンパイラはしばしば同じ効果を達成するためにシフト命令を使用したトリッキーなコードを生成することです。
例えば、Thumbモードで\INS{AND Rdest, Rsrc, 1}または\INS{TST Rsrc, 1}コードを生成するのはトリッキーです。
そのため、Keilは入力\IT{idx}を31ビット左にシフトし、結果の値がゼロかどうかをチェックするコードを生成します。

\subsubsection{セッター}

セッターコードの前半はゲッターと非常に似ており、最初にトリプレットのアドレスが計算され、次に適切なハンドラーのコードにディスパッチするためのジャンプが発生します。

\lstinputlisting[style=customasmARM]{advanced/380_FAT12/Keil_thumb_O3_setter.asm}

\subsection{最適化Keil 5.05（ARMモード）}

\subsubsection{ゲッター}

ARMモード用のゲッター関数には条件分岐がまったくありません！
ARMモードの多くの命令に供給できるサフィックス（\IT{-EQ}、\IT{-NE}など）のおかげで、対応するフラグが設定されている場合にのみ命令が実行されます。

ARMモードの多くの算術命令は、\INS{LSL \#1}のようなシフトサフィックスを持つことができます
（これは、最後のオペランドが1ビット左にシフトされることを意味します）。

\lstinputlisting[style=customasmARM]{advanced/380_FAT12/Keil_ARM_O3_getter.asm}

\subsubsection{セッター}

\lstinputlisting[style=customasmARM]{advanced/380_FAT12/Keil_ARM_O3_setter.asm}

\IT{idx*1.5}の値が2回計算されます。これはKeilコンパイラが生成した冗長性で、除去できます。
アセンブリ関数をより短くするように書き直すこともできます。テストを忘れないでください！

\subsection{（32ビットARM）ThumbとARMモードでのコード密度の比較}

ARM CPUのThumbモードは、ARMモードの32ビット命令の代わりに、より短い（16ビット）命令を作るために導入されました。
しかし、ご覧のとおり、それが価値があったかどうかは言い難いです：ARMモードのコードは常により短いです（ただし、命令はより長いです）。

\subsection{ARM64用最適化GCC 4.9.3}

\subsubsection{ゲッター}

\lstinputlisting[style=customasmARM]{advanced/380_FAT12/ARM64_getter.lst}

ARM64には新しいクールな命令\INS{UBFIZ}（\IT{Unsigned bitfield insert in zero, with zeros to left and right}）があります。
これは、一つのレジスタから別のレジスタに指定された数のビットを配置するために使用できます。
これは別の命令\INS{UBFM}（\IT{Unsigned bitfield move, with zeros to left and right}）のエイリアスです。
\INS{UBFM}は、\INS{LSL/LSR}（ビットシフト）の代わりにARM64で内部的に使用される命令です。

\subsubsection{セッター}

\lstinputlisting[style=customasmARM]{advanced/380_FAT12/ARM64_setter.lst}

\subsection{MIPS用最適化GCC 4.4.5}

ジャンプ/分岐命令の後の各命令が最初に実行されることを覚えておく必要があります。
これはRISC CPUの専門用語で\IT{分岐遅延スロット}と呼ばれます。
物事をシンプルにするために、分岐またはジャンプ命令で始まる各命令ペアで命令を（精神的に）交換するだけです。

MIPSにはフラグがありません（明らかに、データ依存性を簡素化するため）。そのため、分岐命令（\INS{BNE}など）は比較と分岐の両方を行います。

関数プロローグにはGP（グローバルポインタ）設定コードもありますが、これまでのところ無視できます。

\subsubsection{ゲッター}

\lstinputlisting[style=customasmMIPS]{advanced/380_FAT12/MIPS_getter.asm}

\subsubsection{セッター}

\lstinputlisting[style=customasmMIPS]{advanced/380_FAT12/MIPS_setter.asm}

\subsection{実際のFAT12との違い}

実際のFAT12テーブルは少し異なります：\url{https://en.wikipedia.org/wiki/Design_of_the_FAT_file_system\#Cluster_map}。

偶数要素の場合：

\begin{lstlisting}
 +0 +1 +2
|23|.1|..|..
\end{lstlisting}

奇数要素の場合：

\begin{lstlisting}
 +0 +1 +2
|..|3.|12|..
\end{lstlisting}

LinuxカーネルのFAT12関連関数：\\
\href{https://github.com/torvalds/linux/blob/de182468d1bb726198abaab315820542425270b7/fs/fat/fatent.c#L117}{fat12\_ent\_get()}、
\href{https://github.com/torvalds/linux/blob/de182468d1bb726198abaab315820542425270b7/fs/fat/fatent.c#L153}{fat12\_ent\_put()}。

それでも、デモンストレーションのために、バイトレベルのGDBダンプで値がより見やすく認識できるように、このようにしました。

\subsection{演習}

おそらく、ゲッター/セッター関数がより高速になるようにデータを格納する方法があるかもしれません。
値をこのように配置した場合：

（偶数要素）

\begin{lstlisting}
 +0 +1 +2
|23|1.|..|..
\end{lstlisting}

（奇数要素）

\begin{lstlisting}
 +0 +1 +2
|..|.1|23|..
\end{lstlisting}

このデータ格納スキーマは、少なくとも1つのシフト操作を除去することを可能にします。
演習として、私のC/C++コードをそのように書き直し、コンパイラが何を生成するかを見ることができます。

\subsection{まとめ}

ビットシフト（C/C++で\TT{<<}と\TT{>>}、x86で\INS{SHL}/{SHR}/{SAL}/{SAR}、ARMで\INS{LSL}/\INS{LSR}、MIPSで\INS{SLL}/\INS{SRL}）は、ビットを特定の場所に配置するために使用されます。

AND演算（C/C++で\TT{\&}、x86/ARMで\INS{AND}）は、分離中にも不要なビットを削除するために使用されます。

OR演算（C/C++で\TT{|}、x86/ARMで\INS{OR}）は、複数の値を1つにマージまたは結合するために使用されます。
1つの入力値は、別の値が情報を運ぶビットを持つ場所にゼロスペースを持つ必要があります。

ARM64には、あるレジスタから別のレジスタに特定のビットを移動する新しい命令\INS{UBFM}、\INS{UFBIZ}があります。

\subsection{結論}

FAT12は現在どこでもほとんど使用されていませんが、スペース制約があり、12ビットに制限された値を格納する必要がある場合は、FAT12テーブルで行われる方法で密にパックされた配列を検討することができます。